% The cases of test-keep-leadin.tex and test-keep-leadin-off.tex.
\makeatletter
% \roomfor{n}: a new page, filled with a rule so that n lines of text fit under it.
\newcommand\roomfor[1]{%
  \clearpage
  \noindent\rule{1pt}{\dimexpr\textheight-\topskip-#1\baselineskip-\baselineskip\relax}\par}
\newwrite\@leadtest@out
\newcount\@leadtest@i
% \testlisting{name}{n}: a listing of n lines, the first with the label <name>L.
\newcommand\testfile[4]{%
  \immediate\openout\@leadtest@out=\jobname-#1.txt
  \@leadtest@i\z@
  \@whilenum\@leadtest@i<#2\do{\advance\@leadtest@i\@ne
    \immediate\write\@leadtest@out{%
      line \the\@leadtest@i#3\ifnum\@leadtest@i=\@ne\string\label{#1L}\fi#4}}%
  \immediate\closeout\@leadtest@out}
\newcommand\testlisting[2]{%
  \testfile{#1}{#2}{}{}%
  \VerbatimInput[commandchars=\\\{\}]{\jobname-#1.txt}}
% \case{name}: the line of a case, on the line of what it is about, for the log.
\newcommand\case[1]{\typeout{leadin case #1 line \the\inputlineno}}
% \results{name}: the pages of the lead-in, <name>, and the first line, <name>L.
\newcommand\results[1]{\typeout{leadin result #1 lead=\getpagerefnumber{#1}
    first=\getpagerefnumber{#1L}}}
\makeatother

\begin{document}

\roomfor{2}
Lead-in M.\label{m}

\marginpar{note M}
\case{m}\testlisting{m}{3}
Text.

\roomfor{2}
Lead-in N.\label{n}

\ltnote{note N}
\case{n}\testlisting{n}{3}
Text.

\roomfor{2}
Lead-in F.\label{f}

\begin{figure}[b]\rule{1pt}{1pt}\caption{float F}\end{figure}
\case{f}\testlisting{f}{3}
Text.

\roomfor{2}
Lead-in H.\label{h}

\begin{figure}[h]\rule{1pt}{1pt}\caption{float H}\end{figure}
\case{h}\testlisting{h}{3}
Text.

\roomfor{2}
Lead-in S.\label{s}
\vspace{2pt}
\case{s}\testlisting{s}{3}
Text.

\roomfor{2}
Lead-in B, by bigskip.\label{b}
\bigskip
\case{b}\testlisting{b}{3}
Text.

\roomfor{2}
Lead-in R.\label{r}\begin{figure}[b]\rule{1pt}{1pt}\caption{float R in para}\end{figure}
\case{r}\testlisting{r}{3}
Text.

\roomfor{2}
Lead-in U.\label{u}\begin{figure}[h]\rule{1pt}{1pt}\caption{float U in para}\end{figure}
\case{u}\testlisting{u}{3}
Text.

% A listing that breaks: what counts is where it starts.
\roomfor{2}
Lead-in L, a long listing.\label{l}

\marginpar{note L}
\case{l}\testlisting{l}{30}
Text.

% Through minted, whose lines come from a file of its own, over the full width.
\roomfor{2}
Lead-in I, by lstexample*.\label{i}

\marginpar{note I}
\testfile{i}{3}{|}{|}%
\case{i}\lstexample*{text}{\jobname-i.txt}
Text.

\roomfor{2}
Lead-in J, a note and a minted listing.\label{j}

\marginpar{note J}
\testfile{j}{3}{|}{|}%
\case{j}\inputminted[escapeinside=||]{text}{\jobname-j.txt}
Text.

\roomfor{2}
Lead-in P, to code side by side.\label{p}

\marginpar{note P}
\case{p}\begin{didacticpair}{2}
  \rule{1pt}{8\baselineskip}\label{pL} & right
\end{didacticpair}
Text.

\roomfor{2}
Lead-in W, to code side by side, by vspace.\label{w}
\vspace{2pt}
\case{w}\begin{didacticpair}{2}
  \rule{1pt}{8\baselineskip}\label{wL} & right
\end{didacticpair}
Text.

% A chunk of noweb, woven for minted: its header line comes first.
\makeatletter
\ifdefined\nwbegincode
\makeatother
\roomfor{2}
Lead-in K, a code chunk.\label{k}

\marginpar{note K}
\case{k}\nwbegincode{1}\moddef{hello.py}\endmoddef\nwstartdeflinemarkup\nwenddeflinemarkup
\catcode37=14\relax\catcode126=13\relax\catcode36=3\relax\catcode38=4\relax\catcode35=6\relax\catcode94=7\relax\catcode95=8\relax\catcode13=5\relax\partopsep=0pt\relax\begin{minted}[xleftmargin=\codemargin,vspace=0pt,escapeinside=||]{python}
def main():|\label{kL}|
    print("Hello, world!")
\end{minted}
\nwendcode{}\nwbegindocs{2}\nwdocspar
Text.
\nwenddocs{}
\else
\makeatother
\fi

% With a frame, which comes before the first line.
\roomfor{2}
Lead-in Y2, to a listing in a frame.\label{fr}

\marginpar{note Y2}
\testfile{fr}{3}{}{}%
\case{fr}\VerbatimInput[commandchars=\\\{\},frame=lines]{\jobname-fr.txt}
Text.

% Controls: the lead-in stays with the listing.
\roomfor{2}
Lead-in C, a blank line only.\label{c}

\case{c}\testlisting{c}{3}
Text.

\roomfor{2}
Lead-in Y3, to a listing in a frame.\label{fs}
\testfile{fs}{3}{}{}%
\case{fs}\VerbatimInput[commandchars=\\\{\},frame=lines]{\jobname-fs.txt}
Text.

\roomfor{2}
Lead-in A, with a note in the paragraph.\label{a}\marginpar{note A}
\case{a}\testlisting{a}{3}
Text.

\roomfor{2}
Lead-in D, with a footnote.\label{d}\footnote{a footnote}
\case{d}\testlisting{d}{3}
Text.

\roomfor{3}
\section*{Heading Q\label{q}}
\marginpar{note Q}
\case{q}\testlisting{q}{3}
Text.

\roomfor{12}
Lead-in V, to code side by side, in room.\label{v}
\case{v}\begin{didacticpair}{2}
  \rule{1pt}{8\baselineskip}\label{vL} & right
\end{didacticpair}
Text.

\roomfor{12}
Lead-in Q2, a note and code side by side, in room.\label{q2}

\marginpar{note Q2}
\case{qb}\begin{didacticpair}{2}
  \rule{1pt}{8\baselineskip}\label{q2L} & right
\end{didacticpair}
Text.

% A longer paragraph that breaks: its last lines go along.
\roomfor{2}
Lead-in E, of several lines, which breaks within the paragraph, so that at least
its last line goes to the next page with the listing, and we have nothing to warn
about, since the listing still follows the end of its paragraph on one and the
same page, which is all that the option asks.\label{e}
\case{e}\testlisting{e}{3}
Text.

% A listing that fits, after a note: no page break at all.
\clearpage
Lead-in Z, on a page with room.\label{z}

\marginpar{note Z}
\case{z}\testlisting{z}{3}
Text.

% The break is asked for, so no warning, though they are on different pages.
\roomfor{2}
Lead-in G, by goodbreak.\label{g}
\goodbreak
\case{g}\testlisting{g}{3}
Text.

\roomfor{2}
Lead-in X, by pagebreak with an argument.\label{x}
\pagebreak[2]
\case{x}\testlisting{x}{3}
Text.

\roomfor{2}
Lead-in Y, by clearpage.\label{y}
\clearpage
\case{y}\testlisting{y}{3}
Text.

% No lead-in: the end of a list is a place to break, so these are on their own.
\roomfor{2}
\begin{itemize}\item An item.\label{t}\end{itemize}
\case{t}\testlisting{t}{3}
Text.

\roomfor{3}
Listing O, with one after it.\label{o}
\testlisting{o1}{2}
\case{o}\testlisting{o}{3}
Text.

\clearpage
\typeout{leadin results follow}
\results{m}\results{n}\results{f}\results{h}\results{s}\results{b}\results{r}
\results{u}\results{l}\results{i}\results{j}\results{p}\results{w}
\makeatletter
\ifdefined\nwbegincode\results{k}\fi
\results{fr}
\makeatother
\results{c}\results{a}\results{d}\results{q}\results{e}\results{z}\results{g}
\results{v}\results{fs}\typeout{leadin result qb lead=\getpagerefnumber{q2}
    first=\getpagerefnumber{q2L}}
\results{x}\results{y}\results{t}\results{o}
\typeout{leadin results end}
\end{document}
